\documentclass[11pt]{article}
\usepackage[margin=1in]{geometry}
\usepackage[T1]{fontenc}
\usepackage{amsmath,amssymb}

\title{Research Report: Polynomial Complexity and Ridge Regularization in Validation-Based Model Selection}
\author{Agent Laboratory}
\date{}

\begin{document}
\maketitle

\section{Abstract}
Polynomial degree and ridge regularization jointly determine the flexibility of polynomial regression, but their effects on predictive performance require empirical comparison. We describe a preregistered study asking whether additional polynomial complexity improves validation prediction over a linear fit and whether regularization changes that conclusion. The protocol uses a fixed public split containing 80 training examples and 64 validation examples, with polynomial feature construction, intercept handling, and ridge implementation held fixed by a trusted execution host. Eight candidate configurations pair degrees \(d\in\{1,2,4,8\}\) with regularization strengths \(\alpha\in\{0.0,0.1\}\). Candidates are fitted exclusively on training data and selected by minimum validation mean squared error, with exact ties resolved by lower degree and then lower regularization strength. The paired design permits comparisons of regularized and unregularized fits at each degree and, upon completion, an assessment of whether regularization benefits increase with complexity. The supplied execution evidence covers only one distinct configuration, \(d=1\) and \(\alpha=0.0\), despite two reported CPU executions. For this configuration, the host reports training MSE \(0.005661214895460485\) and validation MSE \(0.007958542384134518\), with independent verification reported as valid. This configuration is therefore selected among completed candidates only. The seven remaining candidates have no supplied measurements, so neither the benefit of polynomial complexity nor its interaction with regularization can be estimated. The approximately \(0.00229733\) difference between validation and training MSE is descriptive and does not independently establish overfitting. Validation is reserved for tuning and does not provide an independent test estimate; hidden-test performance and statistical uncertainty are unavailable. The present report establishes a verified linear baseline and an explicit comparison protocol, while leaving the motivating complexity--regularization hypothesis unresolved.

\section{Introduction}
Polynomial regression provides a controlled setting for examining how model complexity and regularization jointly affect prediction. Increasing polynomial degree expands the set of relationships that a model can represent, while ridge regularization penalizes coefficient magnitude and changes the fitted solution within that representation. These choices address different aspects of estimation: degree determines the available features, and regularization determines how strongly their coefficients are constrained. Greater flexibility can be useful when the response contains nonlinear structure, but it can also increase sensitivity to the particular training sample. Consequently, neither increasing degree nor introducing a positive penalty guarantees lower prediction error. This study asks whether additional polynomial complexity improves validation prediction over a linear fit and whether ridge regularization changes that conclusion. The question is relevant to validation-based model selection because complexity and regularization are often chosen together, yet their separate effects require explicit comparisons.

The main difficulty is distinguishing representational flexibility from predictive benefit using a limited sample. A more complex model can fit training observations more closely without improving predictions on observations excluded from fitting. Moreover, the effect of a coefficient penalty depends on the feature representation and its covariance structure; the same numerical regularization strength need not have an equivalent effect across polynomial degrees. The supplied synopsis of \emph{Ridge Regression: Structure, Cross-Validation, and Sketching} discusses squared-error estimation with a quadratic penalty, the dependence of the ridge solution on feature covariance and regularization, and validation-based regularization selection (arXiv 1910.02373v2). This literature motivates the comparison protocol but does not establish which configuration performs best on the present data. Its discussion of validation bias in a high-dimensional asymptotic regime also provides context for interpreting tuned performance, without supplying a quantitative bias estimate for this experiment. We therefore treat increasing benefits from regularization at higher degrees as a hypothesis requiring local measurements.

Our protocol uses a fixed public task with 80 training examples and 64 validation examples. The trusted host controls polynomial feature construction, intercept handling, and ridge implementation, so these components remain fixed across proposed comparisons. The preregistered candidate sequence is \((1,0.0)\), \((1,0.1)\), \((2,0.0)\), \((2,0.1)\), \((4,0.0)\), \((4,0.1)\), \((8,0.0)\), and \((8,0.1)\), where each pair specifies degree and regularization strength. Each candidate is fitted exclusively on training data and evaluated using validation mean squared error. The distinguishing feature is a paired complexity ladder: regularized and unregularized fits are compared at each degree before examining how their difference changes with complexity. For a completed pair, we define the descriptive regularization benefit as
\[
B_d =
\operatorname{MSE}_{\mathrm{val}}(d,0.0)
-
\operatorname{MSE}_{\mathrm{val}}(d,0.1).
\]
A positive value would indicate lower validation error for the regularized fit at that degree. Selection uses the lowest host-reported validation MSE among completed candidates, with exact ties resolved by lower degree and then lower regularization strength. The contributions of the present report are a specified paired comparison protocol, a verified linear baseline, and an explicit account of the evidence needed to complete the motivating comparison.

The available execution evidence covers only the unregularized linear configuration, \(d=1\) and \(\alpha=0.0\). The host reports training MSE \(0.005661214895460485\) and validation MSE \(0.007958542384134518\), with independent verification reported as valid. Although two CPU executions are recorded, the host reports only one distinct evaluated configuration. Thus, the linear configuration is selected among completed candidates only; this selection does not establish superiority over any polynomial or regularized alternative. The validation--training difference is approximately \(0.00229733\), a descriptive observation that does not by itself establish overfitting. No confidence interval, standard error, or significance test is supplied. The remaining seven candidates have no supplied measurements, so neither \(B_d\) nor its dependence on degree can currently be estimated. Completing those evaluations in preregistered order is the proposed next step. Throughout the report, validation error is interpreted as a tuning criterion rather than an independent test estimate. Hidden-test performance is unavailable, and the current conclusion is limited to a verified baseline with an unresolved complexity--regularization hypothesis.

\section{Background}
We consider supervised regression with scalar inputs \(x\in\mathbb{R}\) and responses \(y\in\mathbb{R}\). The supplied observations are partitioned into a training set \(\mathcal{D}_{\mathrm{tr}}=\{(x_i,y_i)\}_{i=1}^{80}\) and a validation set \(\mathcal{D}_{\mathrm{val}}=\{(x_j^{\mathrm{val}},y_j^{\mathrm{val}})\}_{j=1}^{64}\). Training observations determine fitted coefficients, whereas validation observations provide the criterion for selecting model configurations. Polynomial regression represents predictions using powers of the input; in a conventional monomial representation, a degree-\(d\) predictor takes the form
\[
f_d(x)=\beta_0+\sum_{k=1}^{d}\beta_k x^k.
\]
Degree one gives an affine predictor, and higher degrees permit curvature and more complicated input--response relationships. Polynomial regression remains linear in its fitted coefficients even when its prediction is nonlinear in \(x\). In this experiment, the host determines the actual feature construction and intercept handling. The expression above introduces the model family rather than specifying a replacement implementation. The task does not supply a population regression function or establish a particular noise distribution. Therefore, an apparent relationship in the supplied observations cannot be treated as knowledge of the data-generating mechanism.

Least-squares estimation chooses coefficients by minimizing squared residuals on training observations. Ridge regression supplements this fitting criterion with a quadratic coefficient penalty. To describe the mathematical principle without assuming undocumented host conventions, let \(X_d\) denote a degree-specific training design matrix, let \(\mathbf{y}\) denote the training response vector, and let \(P_d\) be a positive semidefinite matrix specifying which coefficients are penalized. A conventional sum-of-squares formulation is
\[
\widehat{\boldsymbol{\beta}}_{d,\lambda}
\in\operatorname*{arg\,min}_{\boldsymbol{\beta}}
\left\{
\|\mathbf{y}-X_d\boldsymbol{\beta}\|_2^2
+\lambda\boldsymbol{\beta}^{\mathsf T}P_d\boldsymbol{\beta}
\right\}.
\]
When \(X_d^{\mathsf T}X_d+\lambda P_d\) is invertible, this formulation gives
\[
\widehat{\boldsymbol{\beta}}_{d,\lambda}
=
\left(X_d^{\mathsf T}X_d+\lambda P_d\right)^{-1}
X_d^{\mathsf T}\mathbf{y}.
\]
These equations are background identities, not a claim that the host uses this precise normalization or matrix representation. For example, penalizing an intercept and leaving it unpenalized correspond to different choices of \(P_d\); using mean squared residuals instead of their sum changes the numerical interpretation of the penalty parameter. The experiment retains the host's conventions and varies only its exposed parameters \(d\) and \(\alpha\). The supplied synopsis of \emph{Ridge Regression: Structure, Cross-Validation, and Sketching} motivates this dependence on feature covariance and regularization strength (arXiv 1910.02373v2).

Polynomial degree and ridge strength affect different components of the prediction problem. Degree changes the representation, while the penalty changes estimation within that representation. For nested polynomial spaces, exact unregularized least-squares minimization cannot increase the minimum training squared error when degree increases: the larger space contains the smaller space's predictors. This mathematical property does not imply that validation error decreases. Additional coefficients can respond to sample-specific fluctuations, and correlated polynomial features can make coefficient estimates sensitive to the observations used for fitting. A positive penalty can reduce such sensitivity while introducing shrinkage bias. Its predictive value consequently depends on the relationship being approximated, the observed design, and the penalty convention. In the simplified case \(P_d=I\), if \(X_d^{\mathsf T}X_d=Q\operatorname{diag}(s_1,\ldots,s_p)Q^{\mathsf T}\) with positive eigenvalues, ridge attenuates the least-squares coefficient component in direction \(k\) by \(s_k/(s_k+\lambda)\). This identity explains why directions with smaller eigenvalues can receive stronger relative shrinkage. It also explains why the same numerical penalty need not have the same effect at different degrees. The identity does not establish that regularization improves prediction on the present split or that its benefit increases with complexity.

Prediction quality is assessed using the validation mean squared error,
\[
V(d,\alpha)
=
\frac{1}{64}\sum_{j=1}^{64}
\left[
y_j^{\mathrm{val}}
-\widehat f_{d,\alpha}(x_j^{\mathrm{val}})
\right]^2,
\]
where \(\widehat f_{d,\alpha}\) is fitted exclusively on \(\mathcal{D}_{\mathrm{tr}}\). The planned candidate set is \(\mathcal{C}=\{1,2,4,8\}\times\{0.0,0.1\}\). For each completed degree pair, the contrast \(B_d=V(d,0.0)-V(d,0.1)\) describes the observed benefit of the positive penalty; a positive value indicates lower validation error at \(\alpha=0.1\). Using identical validation observations makes the contrast paired, but does not by itself supply an uncertainty estimate. Selection minimizes \(V\) over the completed subset of \(\mathcal{C}\), with exact ties resolved by lower degree and then lower \(\alpha\). This rule identifies the best measured candidate under the fixed criterion, not an optimum over all possible polynomial models. Moreover, validation participates in tuning: once it determines a configuration, its error is not an independent test assessment of that selection. The cited ridge synopsis discusses validation bias in a high-dimensional asymptotic setting, but supplies no quantitative correction applicable to this experiment. The current evidence contains only the unregularized linear configuration, so the paired contrasts remain unmeasured. Hidden-test performance, sampling uncertainty, and the complexity--regularization interaction cannot be inferred from the background theory alone.

\section{Related Work}
\emph{Ridge Regression: Structure, Cross-Validation, and Sketching} provides the closest methodological context for this study (arXiv 1910.02373v2). The supplied synopsis describes squared-error estimation with a quadratic coefficient penalty, the dependence of the ridge solution on feature covariance and regularization strength, and regularization selection through held-out or cross-validation error. Our experiment uses the same general estimation principle but asks a narrower empirical question: whether increasing polynomial degree improves validation prediction, and whether a positive ridge penalty changes that comparison. Degree determines the available representation, whereas regularization constrains estimation within that representation. Their effects therefore require separate comparisons. The proposed design pairs \(\alpha=0.0\) and \(\alpha=0.1\) at each degree \(d\in\{1,2,4,8\}\), while keeping the host's feature construction, intercept handling, and fitting implementation fixed. This pairing makes the local question explicit, but it does not introduce a new ridge estimator or establish an unprecedented experimental design. The cited source motivates attention to covariance and regularization; it supplies neither an optimal degree nor an optimal penalty for the present data.

The validation procedure also differs in scope from the theoretical setting described in the ridge paper's synopsis. That synopsis discusses validation bias in a high-dimensional asymptotic regime, whereas our protocol uses a single fixed split with 80 training observations and 64 validation observations. We do not establish that the assumptions of its asymptotic analysis hold here, and we do not transfer a quantitative bias estimate to this experiment. The relevant methodological connection is that validation error guides hyperparameter selection and consequently must be distinguished from an independent assessment of the selected model. Our paired comparison is defined by
\[
B_d=\operatorname{MSE}_{\mathrm{val}}(d,0.0)
    -\operatorname{MSE}_{\mathrm{val}}(d,0.1).
\]
A positive \(B_d\) would describe a validation benefit from regularization at degree \(d\); whether this benefit increases with degree remains a hypothesis. Cross-validation would require a different resampling protocol, while sketching would change the computation being compared. Neither is included in the fixed host experiment. Accordingly, the relevant experimental comparison is between the preregistered ridge configurations, rather than a claimed reproduction of the cited paper's asymptotic or sketching results.

\emph{Agent Laboratory: Using LLM Agents as Research Assistants} concerns the organization of research work rather than polynomial regression itself (arXiv 2501.04227v2). Its supplied summary describes specialized agents performing literature review, planning, experimentation, interpretation, and reporting. This is relevant to the separation between proposing an experiment and obtaining evidence for its outcome. In the present setting, the reporting model receives literature summaries and execution records, while the trusted host performs numerical fitting and validation. The workflow therefore shares the idea of separating research activities, but the supplied experiment does not compare agent architectures or measure research automation performance. Agent Laboratory is not a predictive baseline for this regression task: its described contribution addresses research assistance rather than an alternative mapping from the supplied inputs to responses. Any comparison of research workflows would require additional outcome definitions and execution evidence beyond validation MSE.

\emph{AgentRxiv: Towards Collaborative Autonomous Research} addresses sharing and retrieving prior research reports through a local preprint server and similarity-based retrieval (arXiv 2503.18102v1). Its connection to this study is the use of prior reports as context for planning and interpretation. Our literature context is instead restricted to the supplied frozen corpus; no retrieval strategy or collaborative research system is evaluated. Retrieved conclusions would also remain literature claims rather than measurements on the present split. This distinction limits the conclusions that can be drawn from all three sources. The only supplied local configuration is \(d=1,\alpha=0.0\), with host-reported validation MSE \(0.007958542384134518\). No regularized counterpart or higher-degree result is available, so the proposed paired comparisons cannot yet be reported. The existing literature supports the rationale for validation-based selection and explicit evidence tracking, while the present empirical contribution remains a verified linear baseline and an incomplete comparison protocol.

\section{Methods}
We use a paired comparison protocol to examine polynomial degree and ridge regularization while retaining the fixed training and validation partitions. Let \(\mathcal{D}_{\mathrm{tr}}\) contain the 80 supplied training observations and let \(\mathcal{D}_{\mathrm{val}}\) contain the 64 supplied validation observations. The trusted host determines polynomial feature construction, intercept handling, penalty normalization, and numerical fitting. These implementation choices remain fixed throughout the comparison; only degree \(d\) and regularization strength \(\alpha\) vary. We denote the resulting fitting operation by
\[
\widehat f_{d,\alpha}
=
\mathcal{H}\!\left(\mathcal{D}_{\mathrm{tr}};d,\alpha\right),
\]
where \(\mathcal{H}\) represents the host's polynomial ridge implementation. This notation avoids imposing a coefficient-penalty convention beyond the supplied execution evidence. In particular, we do not assume that the intercept is penalized or that the host's \(\alpha\) equals the parameter in a particular sum-of-squares formulation. All fitted coefficients are determined exclusively from training observations. Validation responses enter only the evaluation and selection procedure. No additional observations, transformations, alternative fitting implementations, or changes to the split are introduced. The methodological motivation is that polynomial degree controls representational flexibility, whereas ridge regularization modifies estimation within that representation. The supplied ridge synopsis supports examining these choices through held-out error and attention to feature covariance (arXiv 1910.02373v2), but does not predict the outcome of the present comparison.

The preregistered candidate set is \(\mathcal{C}=\{1,2,4,8\}\times\{0.0,0.1\}\), with the ordered evaluation sequence
\[
\mathcal{S}
=
\bigl(
(1,0.0),(1,0.1),
(2,0.0),(2,0.1),
(4,0.0),(4,0.1),
(8,0.0),(8,0.1)
\bigr).
\]
Each proposed evaluation consists of one literal configuration assignment containing only the fields \texttt{degree} and \texttt{alpha}; the host performs training and reports the resulting metrics. The sequence is retained regardless of intermediate validation results, subject to available execution resources. This restriction prevents the planned comparison from being replaced by an adaptively expanded search. For each completed candidate, the evaluation criterion is
\[
V(d,\alpha)
=
\frac{1}{64}
\sum_{j=1}^{64}
\left(
y_j^{\mathrm{val}}
-
\widehat f_{d,\alpha}(x_j^{\mathrm{val}})
\right)^2.
\]
The corresponding training diagnostic is
\[
T(d,\alpha)
=
\frac{1}{80}
\sum_{i=1}^{80}
\left(
y_i-\widehat f_{d,\alpha}(x_i)
\right)^2.
\]
Selection uses \(V\), rather than \(T\), because the motivating question concerns prediction on observations excluded from coefficient fitting. Both quantities are reported only when supplied by the host. Training error provides a descriptive account of in-sample fit, but neither a lower training error nor a larger validation--training difference independently establishes a predictive advantage or overfitting. Candidate parameters and planned evaluations are therefore recorded separately from completed measurements.

The paired analysis compares regularized and unregularized predictions on the same validation observations at each degree. For a degree with both candidates completed, define the regularization benefit as
\[
B_d=V(d,0.0)-V(d,0.1).
\]
A positive \(B_d\) indicates lower observed validation MSE for the regularized candidate; a negative value indicates lower error for the unregularized candidate. Polynomial complexity is assessed relative to the degree-one fit at the same regularization strength:
\[
C_{d,\alpha}=V(1,\alpha)-V(d,\alpha),
\qquad d\in\{2,4,8\}.
\]
Thus, positive \(C_{d,\alpha}\) indicates a validation improvement from increasing degree under the specified penalty. To describe whether regularization changes the predictive value of additional complexity, we use
\[
I_d
=
C_{d,0.1}-C_{d,0.0}
=
B_d-B_1.
\]
A positive \(I_d\) would indicate that the regularized comparison favors additional complexity more than the unregularized comparison does. The proposed assessment of increasing regularization benefit examines the signs and magnitudes of \(B_2-B_1\), \(B_4-B_2\), and \(B_8-B_4\). These are discrete contrasts across the selected degrees, not estimates of a continuous relationship. Every contrast requires measurements for all configurations appearing in its definition; missing measurements are not replaced by zeros, theoretical expectations, or extrapolated values. The common validation split makes the comparisons paired, but aggregate MSE values alone do not provide the covariance of per-observation loss differences. Consequently, this protocol reports descriptive contrasts without assigning confidence intervals or statistical significance unsupported by the execution evidence.

Let \(\mathcal{E}\subseteq\mathcal{C}\) denote the distinct configurations with completed, host-reported evaluations available for selection. We choose
\[
(d^\star,\alpha^\star)
=
\operatorname*{arg\,min}_{(d,\alpha)\in\mathcal{E}}^{\mathrm{lex}}
\bigl(V(d,\alpha),d,\alpha\bigr),
\]
where lexicographic ordering first minimizes validation MSE, then resolves exact ties by lower degree and subsequently lower \(\alpha\). Comparisons use the reported metric precision rather than rounded table entries. The selected configuration is described as best among \(\mathcal{E}\), without extending that claim to unmeasured candidates or to all admissible degrees and penalties. The supplied record currently establishes \(\mathcal{E}=\{(1,0.0)\}\): execution \texttt{execution-0002} reports \(T(1,0.0)=0.005661214895460485\) and \(V(1,0.0)=0.007958542384134518\), with independent verification reported as valid. Although the host reports two CPU executions, it reports only one distinct configuration; execution counts are not treated as additional comparison conditions or independent samples. The current selection is therefore \((1,0.0)\), and all regularization and higher-degree contrasts remain unavailable. Completing the remaining seven configurations in the specified order is proposed work, not a completed analysis. Throughout selection, hidden-test data remain excluded. Because validation error determines the selected configuration, its reported value is a tuning criterion and not an independent test estimate of the selected model's generalization performance.

\section{Experimental Setup}
The experiment uses the fixed public regression task \texttt{dev-linear}, task version \texttt{polynomial-ridge-1}, with seed \(7\). Each observation contains a scalar input \(x\), a scalar response \(y\), and a split-specific identifier. The supplied training partition contains \(80\) observations, and the supplied validation partition contains \(64\) observations. These partitions are used exactly as provided: training observations determine fitted coefficients, while validation observations determine the model-selection criterion. The supplied split manifest records generation using \texttt{stdlib random.Random}, seed \(7\), and the generation order training, validation, and test. It also provides separate SHA-256 digests for observation identifiers and split contents. These manifest entries document the supplied data provenance; no regeneration or independent digest computation is performed for this report. Although the manifest lists a withheld test partition containing \(96\) observations, its inputs and responses are unavailable for fitting, selection, or evaluation. The experiment introduces no external data, additional resampling, or altered partitions. Seed \(7\) identifies the supplied task instance rather than a collection of independently sampled experimental repetitions. The task does not provide a population regression function, a response-noise distribution, or application-specific response units. Accordingly, evaluation concerns prediction on this fixed validation partition, and MSE is expressed in squared response units without assigning an application-level interpretation to its magnitude.

The trusted host performs polynomial ridge fitting using its fixed feature construction, intercept handling, and numerical implementation. Candidate submissions expose only \texttt{degree} and \texttt{alpha}; they do not replace the host's fitting code or specify alternative preprocessing. The admissible parameter domain is integer degree \(1\) through \(8\) and finite regularization strength \(0\) through \(100\), but the preregistered experiment considers only degrees \(\{1,2,4,8\}\) and penalties \(\{0.0,0.1\}\). This restricted grid implements the planned paired complexity ladder. Degree one supplies the linear reference, and degrees two, four, and eight provide progressively richer polynomial representations. At each degree, the two penalties permit a comparison between an unregularized fit and a fit with the specified positive penalty. Their numerical values are interpreted under the host's conventions; the available evidence does not establish a particular loss normalization, feature scaling, intercept penalty, or solver tolerance. Those details remain fixed rather than being inferred or modified. The planned evaluation order and currently documented coverage are
\[
\begin{array}{c|c|c|l}
\text{Order} & d & \alpha & \text{Supplied measurement status}\\
\hline
1 & 1 & 0.0 & \text{Host result available}\\
2 & 1 & 0.1 & \text{No supplied measurement}\\
3 & 2 & 0.0 & \text{No supplied measurement}\\
4 & 2 & 0.1 & \text{No supplied measurement}\\
5 & 4 & 0.0 & \text{No supplied measurement}\\
6 & 4 & 0.1 & \text{No supplied measurement}\\
7 & 8 & 0.0 & \text{No supplied measurement}\\
8 & 8 & 0.1 & \text{No supplied measurement}
\end{array}
\]
Each candidate is submitted separately through one literal configuration assignment. The planned order is retained subject to host resources, rather than changed in response to intermediate validation errors. The remaining seven entries describe proposed evaluations and have no reported outcomes.

For a candidate fitted exclusively on training observations, the primary metric is validation mean squared error, and training mean squared error is retained as a diagnostic:
\[
V(d,\alpha)=\frac{1}{64}\sum_{j=1}^{64}
\left[y_j^{\mathrm{val}}-\widehat f_{d,\alpha}
(x_j^{\mathrm{val}})\right]^2,
\qquad
T(d,\alpha)=\frac{1}{80}\sum_{i=1}^{80}
\left[y_i-\widehat f_{d,\alpha}(x_i)\right]^2.
\]
Neither metric includes the ridge penalty term. All candidate comparisons use the same validation observations, ensuring that differences in evaluation sets do not confound the planned degree and penalty comparisons. For each completed degree pair, the descriptive regularization contrast is \(B_d=V(d,0.0)-V(d,0.1)\); a positive value indicates lower validation MSE at the positive penalty. Changes in these contrasts across the preregistered degrees address whether regularization's observed benefit increases with complexity. These comparisons require both measurements at each relevant degree and cannot be calculated from the linear baseline alone. Selection minimizes the host-reported validation MSE among distinct completed candidates, resolving exact ties by lower degree and then lower \(\alpha\). Full supplied metric precision is retained for selection, so presentation rounding does not create artificial ties. Validation serves tuning throughout this procedure and does not become an independent test estimate after a configuration is selected. The setup includes no cross-validation folds, independent split repetitions, or supplied uncertainty estimates; repeated execution on the same fixed observations does not supply those quantities.

The available execution record is identified as \texttt{execution-0002} and classified by the host as \texttt{actual\_cpu\_execution}. Its configuration is \(d=1,\alpha=0.0\), with training MSE \(0.005661214895460485\) and validation MSE \(0.007958542384134518\). The host reports independent verification as \texttt{valid} and test access as \texttt{withheld}. Its artifact list identifies a registered specification, execution log, implementation source files, split manifest, training and validation data files, fitted model, and predictions. These are host-listed provenance artifacts, not files independently inspected for this report. The host additionally reports two actual CPU attempts to date, one distinct evaluated configuration, an execution limit of \(51\), and \(49\) remaining CPU executions at the recorded point. Attempt counts and distinct configuration counts describe different aspects of execution coverage; neither implies that the remaining candidates were evaluated. Although the stated remaining CPU allowance exceeds the seven pending candidates, resource availability is not evidence of completed measurements, and the record does not establish why the search remains incomplete. Hardware specifications, elapsed execution times, and software dependency versions are not supplied. The present setup therefore supports reporting the verified baseline and selecting \((1,0.0)\) within the completed singleton candidate set, while reserving all paired comparisons and higher-degree conclusions for additional trusted-host results.

\section{Results}
The supplied execution evidence establishes one measured configuration: the unregularized linear model with \(d=1\) and \(\alpha=0.0\). Fitted exclusively on the \(80\) training observations, this configuration achieved training MSE \(0.005661214895460485\) and validation MSE \(0.007958542384134518\) on the \(64\) validation observations. These values are reported directly from the trusted host's \texttt{execution-0002} record, which identifies an actual CPU execution and reports independent verification as \texttt{valid}. Table~\ref{tab:results-coverage} distinguishes the measured baseline from the seven preregistered configurations without supplied results. Missing entries represent unavailable measurements, rather than zero error, failed fitting, or evidence of inferior performance. The host lists model and prediction artifacts, but their contents were not independently inspected for this report. Consequently, the analysis uses the supplied aggregate metrics and does not reconstruct coefficients, residual distributions, or additional performance measures.
\begin{table}[ht]
\centering
\begin{tabular}{cccl}
\hline
Degree & \(\alpha\) & Validation MSE & Evidence status\\
\hline
1 & 0.0 & \(0.007958542384134518\) & Host measurement\\
1 & 0.1 & --- & No supplied measurement\\
2 & 0.0 & --- & No supplied measurement\\
2 & 0.1 & --- & No supplied measurement\\
4 & 0.0 & --- & No supplied measurement\\
4 & 0.1 & --- & No supplied measurement\\
8 & 0.0 & --- & No supplied measurement\\
8 & 0.1 & --- & No supplied measurement\\
\hline
\end{tabular}
\caption{Evidence coverage in preregistered order. Dashes denote unavailable validation results.}
\label{tab:results-coverage}
\end{table}

Applying the specified selection rule to the completed candidate set gives
\[
\mathcal{E}=\{(1,0.0)\},
\qquad
(d^\star,\alpha^\star)=(1,0.0).
\]
This selection is restricted to the single distinct evaluated configuration. No competing configuration has a supplied validation score, so selection does not demonstrate that linear regression outperforms a regularized linear model or a higher-degree polynomial. The exact-tie rule, favoring lower degree and then lower \(\alpha\), has no effect within this singleton set. The host reports two actual CPU attempts but only one distinct configuration. These counts do not establish two independently sampled trials or supply a second numerical result for comparison. Although the record reports \(49\) remaining CPU executions under a limit of \(51\), available capacity does not establish completion of any pending candidate. The supplied evidence does not identify the reason for the incomplete search.

The validation MSE exceeds the training MSE by the descriptive difference
\[
V(1,0.0)-T(1,0.0)
=
0.007958542384134518-0.005661214895460485
\approx 0.00229733.
\]
This calculation summarizes the reported metrics; it is not an additional host measurement or a statistical test. Training and validation errors summarize different observations, and the training observations also determined the fitted coefficients. The difference therefore cannot, by itself, establish overfitting or identify its mechanism. No standard error, confidence interval, significance test, or distribution of observation-level losses is supplied. Aggregate MSE alone does not determine the variability needed for such uncertainty estimates. Repeating an execution on the same fixed split also does not establish variability across independently sampled datasets. The evidence supports reporting the observed errors and their difference, without claiming statistical significance or a population-level prediction guarantee.

None of the planned paired ablations can currently be evaluated. In particular, \(B_d=V(d,0.0)-V(d,0.1)\) remains unavailable at every preregistered degree, including degree one. Higher-degree comparisons against the linear reference and interaction contrasts \(I_d=B_d-B_1\) likewise lack the required measurements. The fixed splits and host implementation define consistent conditions for the proposed comparisons, but cannot substitute for their execution. Completion would require the remaining seven candidates to be evaluated separately in preregistered order, retaining identical fitting and evaluation conventions. Such completion is proposed work, with no predicted outcome assigned here. Validation remains the tuning criterion and is not an independent test estimate. Test access is reported as withheld, and no hidden-test result, timing comparison, monetary cost, or baseline improvement is available. The empirical conclusion is therefore limited to a verified linear baseline; the predictive effects of polynomial complexity, ridge regularization, and their interaction remain unresolved.

\section{Discussion}
The study specifies a paired comparison of polynomial complexity and ridge regularization using 80 training observations and 64 validation observations. The available evidence establishes only the unregularized linear baseline, \(d=1,\alpha=0.0\), with host-reported training MSE \(0.005661214895460485\) and validation MSE \(0.007958542384134518\). Independent verification was reported as valid. This configuration is selected among completed candidates only, because the host reports one distinct evaluated configuration despite two CPU executions. The incomplete search provides no comparative evidence favoring linear prediction over higher-degree or regularized alternatives. Its contribution is a documented baseline and a preregistered protocol for resolving that comparison.

The validation--training difference of approximately \(0.00229733\) is descriptive and does not independently establish overfitting. No sampling uncertainty or significance assessment is available, and repeated executions on the same observations do not constitute independent dataset repetitions. Validation is the selection criterion and cannot also serve as an independent test estimate of the selected configuration. Hidden-test performance remains unavailable. These limitations distinguish numerical verification of an execution from evidence about generalization. Furthermore, a single split and two planned penalty values restrict the scope of any eventual conclusion: completing the grid would identify the lowest observed validation error among those eight candidates, rather than an optimum throughout the admissible parameter domain.

The supplied ridge synopsis motivates examining regularization through feature covariance and held-out prediction error (arXiv 1910.02373v2). It does not determine the direction or magnitude of the local regularization effect. The immediate proposed continuation is to evaluate the remaining seven configurations separately in preregistered order, retaining the host's feature construction, intercept handling, implementation, and fixed splits. Completed measurements would permit calculation of \(B_d=V(d,0.0)-V(d,0.1)\) and \(I_d=B_d-B_1\). Positive values would describe, respectively, a validation benefit from regularization and a larger benefit than at degree one. Neither contrast is currently estimable. Their eventual interpretation should retain the distinction between descriptive paired comparisons and statistical evidence, without attributing an observed pattern to a covariance mechanism that has not been measured.

\end{document}